\section{Animacoes em GIF e incorporacao no PDF}

Os GIFs originais em \codefile{output/gifs} foram convertidos para sequencias de frames PNG dentro de \codefile{relatorio/frames}, usando \codefile{ffmpeg}. Nesta revisao, a animacao incorporada ao PDF foi atualizada para o caso Pi com retorno aterrado: a fonte de surto fica entre o no de referencia e a entrada da bobina, e o terminal final da bobina retorna diretamente ao mesmo no de referencia. Assim, a animacao documenta a condicao solicitada de fim da bobina conectado ao no comum com a fonte, em vez de enfatizar apenas a extremidade aberta.

A incorporacao usa o pacote \codefile{animate}. Em visualizadores compativeis, como Adobe Acrobat Reader, as figuras podem ser reproduzidas com controles. Em visualizadores simples, o PDF pode exibir apenas um quadro estatico.

\subsection{Onda viajante com retorno ao no comum}

\begin{figure}[H]
\centering
\animategraphics[controls,loop,autoplay,width=0.90\linewidth]{5}{frames/voltage_wave_grounded/frame-}{1}{75}
\caption{Animacao da distribuicao de tensao ao longo da bobina no modelo Pi com terminal final conectado ao no de referencia. A entrada recebe o surto e a extremidade em 100\% permanece presa ao potencial comum da fonte.}
\label{fig:anim-grounded-wave}
\end{figure}

A Figura~\ref{fig:anim-grounded-wave} mostra que a resposta deixa de ser a mesma do terminal aberto. No caso aberto, a corrente terminal nula favorece reflexao de tensao com mesmo sinal e elevacao da tensao na extremidade. No caso com retorno ao no comum, o terminal final e imposto em \SI{0}{V}; por isso a informacao mais importante passa a ser a distribuicao interna de tensao ao longo do enrolamento.

\subsection{Mapa de calor animado do caso aterrado}

\begin{figure}[H]
\centering
\animategraphics[controls,loop,autoplay,width=0.94\linewidth]{5}{frames/heatmap_grounded/frame-}{1}{75}
\caption{Mapa de calor animado do modelo Pi com terminal final conectado ao no comum com a fonte. O cursor temporal relaciona a distribuicao instantanea com o historico posicao-tempo.}
\label{fig:anim-grounded-heatmap}
\end{figure}

A Figura~\ref{fig:anim-grounded-heatmap} complementa a animacao da onda viajante porque mostra o comportamento no plano tempo--posicao. A borda em 100\% permanece vinculada ao no de referencia, enquanto os nos internos respondem ao impulso aplicado na entrada. Essa visualizacao torna mais direta a comparacao conceitual com o caso aberto: muda-se a condicao de contorno, mudam-se a reflexao e a distribuicao de solicitacao dieletrica.
